\documentclass[11pt]{article}
\usepackage[T1]{fontenc}
\usepackage{lmodern,microtype,amsmath,amssymb,amsthm,booktabs}
\usepackage[margin=27mm]{geometry}
\usepackage[hidelinks]{hyperref}
\usepackage{enumitem}
\setlist{nosep}
\newtheorem{theorem}{Theorem}
\newtheorem{lemma}[theorem]{Lemma}
\newtheorem{corollary}[theorem]{Corollary}
\newcommand{\ind}{\mathbf 1}
\newcommand{\E}{\mathbb E}
\newcommand{\card}[1]{\lvert #1\rvert}
\newcommand{\floor}[1]{\lfloor #1\rfloor}
\newcommand{\ceil}[1]{\lceil #1\rceil}
\title{A Root-Separation Barrier for Tensor-Interpolation SAT Algorithms}
\author{}
\date{26 November 2026}
\begin{document}
\maketitle
\begin{abstract}
Consider a fixed signed sum of local products that recovers the model count
of a pinned cubic expander modulo a nontrivial residue ring. At each term the
coordinate, local functions and coefficient may differ. We prove a lower
bound of \(2^{\Omega_\gamma(N/\sqrt{\log N})}\) on the number of terms under a
collision certificate on the port labels, where \(N\) is the vertex count
and \(\gamma>0\) is the edge expansion. For products over all symbol groups
the bound is uniform over alphabet sizes \(r\ge 27\) and requires only a
nonempty set of retained coordinates. A second result allows a different
balanced union of symbols per term, with constants depending on the alphabet
and balance. The proof selects a single oriented induced matching on which
every eligible term has a rank deficit. The resulting cost bound applies to
traces that materialize every term. We give the finite inequalities and the
scope of their Lean proofs.
\end{abstract}

\section{The output and the layout}

The fixed-density counting algorithm of OpenAI~\cite{source} motivates a
question about the expression at its root. If the same output admits many
decompositions into products of group-local functions, how short can such
a decomposition be? We study this question as an extremal problem with
an explicit output contract.

Balanced products, induced matchings and matrix rank also appear in the
structured arithmetic-circuit lower bounds of de Colnet and Mengel~\cite{dm}.
Here the terms may choose different coordinate cuts, so the rank estimate
requires one matching and orientation that work for all eligible cuts.

Take a nonempty finite vertex set \(V\), \(\card V=N\), and a loopless cubic
graph with a fixed orientation of its edges \(E\). Assume that for every
\(S\subseteq V\),
\[
   |\partial S|\ \ge\ \gamma\min\{|S|,N-|S|\},\qquad 0<\gamma\le 1,
\]
with rational \(\gamma\). For each sign vector \(\sigma\in\{0,1\}^V\), impose
one unit clause \(x_v=\sigma_v\) per vertex and one positive clause
\(x_v\lor x_w\) per edge. The model count is the Boolean number
\[
    C(\sigma)=\prod_{(v,w)\in E}(\sigma_v\lor\sigma_w).
\]
All signs vary independently. We work on each supplied expanding graph.
Infinite families of cubic bipartite expanders provide familiar examples
of this hypothesis~\cite{mss}.

Let \(A\) be a finite nonempty alphabet, \(r=|A|\), \(I\) the full coordinate
set, and \(I_0\subseteq I\) the retained coordinates. A layout consists of
\(M=10N\) ports, an injection from \(V\) into the ports, and a map from ports
to pairs of labels in \(A^I\times A^I\). Write \(L(v)\) for the first label
of the port of \(v\), restricted to \(I_0\). If \(X_c\) counts ports in cell
\(c\) and \(Q=r^{2|I|}\), the certificate is
\begin{equation}\label{eq:certificate}
  \sum_c(X_c-M/Q)^2\le 2M\Delta,\qquad 2Q^2\Delta\le M .
\end{equation}
The factor ten has a concrete origin. In the motivating reduction a variable
appears four times. A three-valent equality tree contributes one internal edge
to its four clause incidences. Subdividing these edges gives \(10N\) wires.
The theorems below take the layout and certificate as inputs.

Let \(R,S\) be commutative rings, \(S\) nontrivial, and \(K\) a field. Fix an
injective unital homomorphism \(e:R\to K\) and a unital homomorphism
\(\pi:R\to S\). The output is \(B:\{0,1\}^V\to R\), with
\(\pi B(\sigma)=C(\sigma)\) in \(S\) for every \(\sigma\). This formulation
allows cancellation of nonintegral summands in \(K\).
In a full-group representation each term \(z\) has coefficient \(c_z\in K\),
one coordinate \(i_z\in I_0\), and one arbitrary \(K\)-valued function per
symbol. Thus
\begin{equation}\label{eq:full}
 eB(\sigma)=\sum_{z\in T}c_z\prod_{a\in A}
       F_{z,a}\bigl(\sigma|_{\{v:L(v)_{i_z}=a\}}\bigr).
\end{equation}
The index set, coordinates, coefficients and functions are fixed for the
whole input family. Let \(h=|T|\).

\section{Finite lower bounds}

Set \(\lambda=\floor{\log_2 N}\),
\(t(\gamma,\eta)=2^{-36}\gamma^2\eta^2/32\), and
\[
 \mathcal C(B,a)=2^{\max\{10,a,4B+1\}},\qquad
 \mathcal M(k,b,t)=\ceil{7776kb/t^2}.
\]
These large cutoffs make the finite statements uniform and keep the
integer rounding explicit.

\begin{theorem}[Full-group representation]\label{thm:full}
Assume \eqref{eq:certificate}, \(r\ge27\), and \(I_0\ne\varnothing\).
Put \(t=t(\gamma,1/4)=2^{-45}\gamma^2\).
If \(N\ge\mathcal C(\mathcal M(1,80,t),0)\), every representation
\eqref{eq:full} admits an integer \(d\ge0\) such that
\begin{equation}\label{eq:main}
      (tN)^2\le 320\lambda d^2,\qquad 2^d\le h .
\end{equation}
In particular \(\log_2 h\ge tN/\sqrt{320\lambda}\).
\end{theorem}

There is a second representation class. A menu entry consists of
\(i\in I_0\) and \(U\subseteq A\), giving the cut
\(P=\{v:L(v)_i\in U\}\). The nonempty menu contains distinct pairs
\((i,U)\). Assume both sides of every cut contain at least \(\eta N\)
vertices, with rational \(0<\eta\le1/2\). A term assigned to this entry is
the product of arbitrary field-valued functions of the two sides, with a
scalar coefficient absorbed into one function. Each term may choose
its own entry. Balance counts independently variable pins.

\begin{theorem}[Balanced menu]\label{thm:menu}
Assume the same graph, output contract and certificate. Put
\(s=\floor{7N/2}\), \(u=\floor{\ceil{\log_2 s}/131072}\),
\(b_r=80\cdot2^{2r}\) and \(t=t(\gamma,\eta)\). Suppose
\(|I_0|\le 2u\) and \(N\ge\mathcal C(\mathcal M(2^r,b_r,t),3)\).
For a sum of \(h\) menu terms there exists an integer \(d\ge0\) with
\[
      (tN)^2\le 4 b_r(\lambda+3)d^2,\qquad 2^d\le h .
\]
For fixed \(r,\eta,\gamma\), this gives the stated exponential order as
a lower bound for any eligible menu, including the menu of all eligible
balanced symbol unions.
\end{theorem}

\begin{corollary}[Complete materialization]\label{cor:cost}
Under the hypotheses of Theorem~\ref{thm:full}, consider a finite trace
whose events are either tagged by one term \(z\in T\) or untagged. Every
term occurs at least once and every event appearing in the trace has a
positive integer charge. Its total charge is at least \(2^d\) for a \(d\)
satisfying \eqref{eq:main}. For every integer \(p\ge0\), both \(h\) and this
charge exceed \(N^p\) once also
\[
     N\ge\mathcal C\bigl(\ceil{320(p+1)^2/t^2},1\bigr).
\]
\end{corollary}

\section{A common rank test}

We give the counting and rank arguments for both theorems. The accompanying
Lean development proves the finite estimates with the constants above.

\paragraph{Density of the actual pins.}
Every cell has \(X_c\le M/Q+\sqrt{2M\Delta}\le2M/Q\) by
\eqref{eq:certificate}. Summing over the unused first-label coordinates and
over the second label shows that the law \(\mu\) of \(L(v)\) for a uniform
vertex is at most \(20\) times the uniform product law on \(A^{I_0}\).
Indeed the injection selects \(N\) ports out of \(10N\), and every
projected label has the same number of full-cell preimages.
In particular a single symbol owns at most \(20N/r\) pins. For \(r\ge27\)
this is at most \(3N/4\). If one group has at least \(N/4\) pins it supplies
a balanced side; otherwise accumulate groups until their total first reaches
\(N/4\), at which point it is less than \(N/2\).
Fix one such union per coordinate in \eqref{eq:full} and multiply the
factors on each side. An occupied cell also gives \(1\le2M/Q\).
Since \(I_0\ne\varnothing\) and \(r\ge27\), we have
\(20\cdot2^{|I|}\le r^{2|I|}\le20N\), hence \(|I_0|\le\lambda\).

\paragraph{A simultaneous imbalance.}
Write \(k=|I_0|\), \(J\) for the menu, and \(R_i\) for its number of
entries at coordinate \(i\). Give its centered cut indicators arbitrary
signs. Under the uniform product law the variance of the signed sum is at
most \(\sum_i R_i^2\), by independence across coordinates. Centering first
at the uniform mean and using density domination bounds the variance on
actual vertices by \(20\sum_i R_i^2\). A common integer scale \(H\ge1\)
can be taken with \(H^2\le80k\) in the one-cut-per-coordinate case and
\(H^2\le b_r k\) for a distinct menu, since \(R_i\le 2^r\).

For each cut take the difference of its indicator between the two endpoints
of an oriented edge. Expansion gives an absolute first moment at least
\(2\gamma\eta/3\) under the uniform edge law. Cubic regularity and
\((x-y)^2\le2x^2+2y^2\) bound the second moment of every signed sum of
these differences by \(4H^2\).
The simultaneous correlation step is a finite form of the sign-maximization
argument of Nazarov~\cite{nazarov}. For clarity, its Huber potential is
quadratic on \([-4,4]\) and extends by its tangent lines outside. Choose
signs maximizing its expected value on the signed sum divided by \(2H/p_0\), where \(p_0=2\gamma\eta/3\).
The derivative divided by four gives edge weights in \([-1,1]\). Comparing
the maximizer with each single-sign flip gives a correlation for every
cut. To see the quantitative estimate, let \(f\) be the edge difference
for that cut and \(w\) the scaled maximizing sum. Then
\(\mathbb E f^2\ge p_0\), \(|f|\le1\), and \(\mathbb E w^2\le p_0\).
The event \(|w|>2\) has probability at most \(p_0/4\), leaving central
energy \(\mathbb E[f^2\ind_{\{|w|\le2\}}]\ge3p_0/4\).
A single flip moves \(w\) by at most two. On the central event the whole
move lies in the quadratic part, and elsewhere the supporting tangent
still gives a lower bound. For \(b\) equal to the derivative divided by four,
maximality gives \(|\mathbb E[b f]|\ge 3p_0^2/(32H)\). As there are
\(3N/2\) edges, the unnormalized correlation is at least
\(\gamma^2\eta^2 N/(16H)\) for every cut at once.

Mark edges independently and keep a marked edge only when its entire
conflict neighborhood contains no other marked edge. Conflicts include
both a shared endpoint and an edge between endpoints. The retained edges
therefore form an induced matching. Thin the orientations with independent
biased signs calibrated to each edge's own retention probability. The
expected signed contribution of edge \(e\) is \(2^{-36}b_e\): if its
retention probability is \(\pi_e\), use an independent orientation of mean
\(2^{-36}b_e/\pi_e\). The cubic conflict block has size at most 36, so
\(\pi_e\ge2^{-36}\). The absolute expected imbalance for each cut is at
least \(2tN/H\). Disjoint mark neighborhoods give independent contributions;
counting potentially dependent ordered pairs bounds the variance by
\(4|E|\,36^2=7776N\). Chebyshev's inequality and a union bound
yield one matching and one
orientation for which the endpoint imbalance of every cut is at least
\(\floor{tN/H}\), provided
\[
                 7776 |J| H^2/t^2<N .
\]
The formal derivation is in \texttt{CommonMatchingSelection},
\texttt{ExpanderBarrier} and \texttt{FiniteGrowth}. The same size condition
ensures \(tN/H\ge2\). Thus rounding reduces the margin by at most a factor
of two. For Theorem~\ref{thm:full}, \(|J|=k\le\lambda\). For
Theorem~\ref{thm:menu}, \(|J|\le2^r k\) and \(k\le\lambda+3\).
The displayed cutoffs ensure the required strict inequality.

\paragraph{Rank on the selected matching.}
Fix all pins outside the endpoints to one. Separate the two endpoint
classes of the oriented matching into row and column bits. If the matching
has \(m\) edges, the matrix of the target count is, up to ordering,
the tensor power of the OR matrix
\[
             \begin{pmatrix}0&1\\1&1\end{pmatrix}^{\otimes m}.
\]
Its determinant is a unit in every nontrivial residue ring. The determinant
of the corresponding matrix of \(B\) therefore has nonzero image under
\(\pi\); injectivity of \(e\) transfers nonvanishing to \(K\). Its field
rank is \(2^m\). On a given cut, let \(a\) and \(b\) count the row and
column endpoints on its first side. A separated product has rank at most
\[
 2^{\min(a,b)+\min(m-a,m-b)}=2^{m-|a-b|}.
\]
The two side matrices act on disjoint row and column coordinates, so their
ranks multiply in this upper bound. Every menu term is tested on the
\emph{same} matrix. Rank subadditivity now gives \(2^m\le h2^{m-d}\),
where \(d=\floor{tN/H}\). The bounds on \(H\) and rounding give the
two theorems. Each distinct term index requires an event in a covering
trace, which gives the cost assertion. The polynomial comparison with its
explicit cutoff is proved in \texttt{PowerCost}.

\section{What the model measures}

A simple computation on the same inputs illustrates the role of the output
contract. Count the edges violated by the pins,
\[
                 W(\sigma)=\sum_{(v,w)\in E}(1-\sigma_v)(1-\sigma_w).
\]
An edge term separates across any fixed cut, including when both endpoints
lie on the same side. The count \(C(\sigma)\) equals the indicator of
\(W(\sigma)=0\). Computing that indicator uses a global zero test after the
sum. This example is formalized in \texttt{PinnedCNF}. Applying the cost
theorem to a specific SAT program would require a correspondence between
its operations and the charged events of Corollary~\ref{cor:cost}.
The verified results here concern representation length and this stated
execution model.

\paragraph{Proof sources.}
The theorem index and build records accompany this paper. The entry points
are \texttt{RetunedLayout.modular\_output\_barrier},
\texttt{MenuBounds.source\_distinct\_menu\_barrier}, and the two
\texttt{RetunedLayout} theorems on superpolynomial representation and cost.
The project uses Lean 4.24.0 and a pinned Mathlib checkout. The research
question and release were directed by the project initiator. AI assistants
contributed to the mathematical development, review, formal proofs and
preparation of this manuscript.

\begin{thebibliography}{9}
\bibitem{source} OpenAI. \emph{The Exponential Time Hypothesis is False}.
25 October 2026 revision, quantitative reduction and contraction.
\href{https://github.com/openai/math/tree/main/preprints/The-Exponential-Time-Hypothesis-is-False-October-25-2026}{Source and manuscript}.
\bibitem{nazarov} F. Nazarov. \emph{The ``Bang solution'' of the
coefficient problem}. Part 2, 1997.
\href{https://www.mathnet.ru/eng/aa774}{Publication record}.
\bibitem{mss} A. W. Marcus, D. A. Spielman and N. Srivastava.
\emph{Interlacing families I: Bipartite Ramanujan graphs of all degrees}.
Annals of Mathematics 182 (2015), 307--325.
\href{https://doi.org/10.4007/annals.2015.182.1.7}{doi:10.4007/annals.2015.182.1.7}.
\bibitem{dm} A. de Colnet and S. Mengel.
\emph{A Compilation of Succinctness Results for Arithmetic Circuits}.
\href{https://arxiv.org/abs/2110.13014}{arXiv:2110.13014}.
\end{thebibliography}
\end{document}
